%! TeX root = ./main.tex
% vim: set spelllang=en:
\chapter*{Conclusion}
\addcontentsline{toc}{chapter}{Conclusion}
\markboth{\textsc{conclusion}}{}

Phylogenetic reconciliation makes it possible to study, in a comparative way, the evolution of units relative to that of their hosts, for example that of gene families compared to that of species.
Among existing reconciliation methods, most assume an independent evolution of the different units that may coexist in a host.
This assumption is invalidated notably for some genes whose organization in the genome makes their functions, and therefore their evolution, interdependent.
One of the methods that make it possible to integrate the interdependence between units into reconciliation is segmental reconciliation.

In this thesis, we have shown that the segmental reconciliation model can be successfully extended to take into account the possibility of horizontal transfers between hosts.
Adding this type of event to the segmental reconciliation model notably makes it possible to consider applications to the study of the evolution of bacterial systems, as we illustrated through applications to CRISPR-Cas systems.
Although the algorithmic complexity of the developed method remains reasonable, this extension of the model revealed the existence of a very large number of co-optimal solutions that hinder the interpretability of the solutions produced.

In order to carry out a systematic study of this extended solution space, we adopted an algebraic approach to the dynamic programming techniques used for the design of the algorithms.
To this end, we presented a unified theory of algebraic dynamic programming showing how cross-cutting questions concerning combinatorial optimization problems can be addressed generically.
This theory was put into practice on the segmental reconciliation problem, notably giving us a glimpse of the potential for an accelerated development of new algorithms in the future through the combination of existing algorithms.

Access to these algebraic tools should make it possible to carry out a more thorough validation of the models developed in this thesis.
Future work will thus have to be carried out to measure what proportion of real scenarios can be recovered by the model and to validate whether segmental reconciliation indeed yields qualitatively better scenarios than bilateral reconciliation or adjacency reconciliation and, if so, under what circumstances.
This work could notably build on existing simulators that make it possible to generate synthetic data.
This validation will make it possible to draw more reliable biological conclusions about the evolution of dependent systems.

Several extensions can be proposed for the methods developed in this manuscript.
The evolution of groups or syntenies is not necessarily tree-like, since they can split and merge over time.
Thus, segmental reconciliation could be extended to phylogenetic networks, that is, to phylogenetic structures containing undirected cycles.
Work has already been carried out to extend bilateral reconciliation to networks and could thus serve as a starting point \parencite{Chan2024EfficientAlgorithmReconciliation}.
Moreover, the developed methods ignore the potential evolution of the linear arrangement of groups of units, for example when these are arranged along a genome.
The algorithmic complexity of the problem associated with ordered segmental reconciliation remains open as far as we know.
Another potential avenue could be to search \emph{a posteriori} for orders on the groups that are compatible with a fixed segmental reconciliation.

Finally, several research directions can be considered in connection with algebraic dynamic programming.
The notion of grammar product considered in \cref{chap:adp-rec} remains to be developed, and the first results obtained by \textcite{Siederdissen2015ProductGrammarsAlignment} in this area remain limited to certain forms of grammars.
Obtaining a general product operator could enable algorithmic advances, for example with regard to multilateral reconciliation~\parencite{Menet2022PhylogeneticReconciliation}.
The development of methods for specifying algorithms on general graphs, for example through their tree decomposition, could also enable the development of algorithms for segmental reconciliation extended to networks.
